A rainha morreu.

Por ela ter sido uma pessoa muito relevante para a história de seu país, o funeral
da rainha será aberto ao público e televisionado ao vivo para toda a nação.
Uma procissão aberta acontecerá durante o funeral, levando o caixão com o corpo
da monarca. A procissão sairá da Igreja Imperial (onde o corpo será velado) e
irá, em linha reta, até o Cemitério Real (onde o corpo será enterrado).

Preocupada com a segurança do evento, a organização da procissão decidiu que a
mesma não andará em velocidade constante durante todo o trajeto. Assim,
a organização demarcou $N$ pontos no trajeto, sendo o primeiro a $x_1$ metros da
Igreja; o segundo a $x_2$ metros da Igreja; etc., até o último ponto, a $x_N$
metros da Igreja, onde fica o Cemitério.

\begin{center}
    \includegraphics[scale=0.8]{funeral/funeral.png}
\end{center}

Então, a
procissão sairá da Igreja, com velocidade inicial 0 (zero) e aceleração constante
$a_0$ m/s$^2$, até o ponto $x_1$. Neste ponto,
a aceleração é alterada para $a_1$ m/s$^2$, com a qual a procissão vai
até o ponto $x_2$. Neste ponto, a aceleração é alterada para $a_2$ m/s$^2$ e a
procissão segue até o ponto $x_3$; e assim por diante, até que a procissão chegue
ao Cemitério.

Para ajudar a organização com o restante do funeral, sua tarefa é determinar em
que horário a procissão chegará no Cemitério Real.


\subsection*{Entrada}

A primeira linha contém dois inteiros $N$ e $a_0$ ($1 \leq N \leq 10^4$, $0 <
        a_0 \leq 100$).
As próximas $N-1$ linhas contém, cada uma, os inteiros $x_i$ e $a_i$
($0 < x_i \leq 10^9$, $-100 \leq a_i \leq 100$).
A última linha contém o inteiro $x_N$ ($0 < x_N \leq 10^9$).

Os valores de $x_i$ são dados em ordem estritamente crescente, iniciando por $x_1$.
É garantido que a velocidade da procissão é sempre positiva durante todo o
trajeto (inclusive em todos os trechos em que a aceleração é negativa), e que é igual a zero apenas no ponto inicial (na Igreja).

\subsection*{Saída}

Imprima uma linha com um valor $T$, indicando que a procissão chegará no
Cemitério Real $T$ segundos após a saída da Igreja Imperial. Arredonde e imprima a resposta
com duas casas decimais.

\newpage
\begin{table}[!h]
\centering
\begin{tabular}{|l|l|}
\hline
\begin{minipage}[t]{3in}
\textbf{Exemplo de entrada}
\begin{verbatim}
1 3
100
\end{verbatim}
\vspace{1mm}
\end{minipage}
&
\begin{minipage}[t]{3in}
\textbf{Exemplo de saída}
\begin{verbatim}
8.16
\end{verbatim}
\vspace{1mm}
\end{minipage} \\
\hline
\end{tabular}
\end{table}

\begin{table}[!h]
\centering
\begin{tabular}{|l|l|}
\hline
\begin{minipage}[t]{3in}
\textbf{Exemplo de entrada}
\begin{verbatim}
2 3
50 5
100
\end{verbatim}
\vspace{1mm}
\end{minipage}
&
\begin{minipage}[t]{3in}
\textbf{Exemplo de saída}
\begin{verbatim}
7.97
\end{verbatim}
\vspace{1mm}
\end{minipage} \\
\hline
\end{tabular}
\end{table}

\begin{table}[!h]
\centering
\begin{tabular}{|l|l|}
\hline
\begin{minipage}[t]{3in}
\textbf{Exemplo de entrada}
\begin{verbatim}
5 10
60 5
120 -5
180 0
200 -6
295
\end{verbatim}
\vspace{1mm}
\end{minipage}
&
\begin{minipage}[t]{3in}
\textbf{Exemplo de saída}
\begin{verbatim}
11.64
\end{verbatim}
\vspace{1mm}
\end{minipage} \\
\hline
\end{tabular}
\end{table}
